% appendix.tex — Specification tests, robustness, and supplementary material

\appendix

\section{Specification Tests}\label{app:spec-tests}

\subsection{Test 1: Sign Restriction}

\textbf{Hypothesis 1.} $\beta_2 > 0$: the quadratic term is positive,
indicating that extreme concentration drives exits.

\textbf{Result: PASS.} $\hat{\beta}_2 = +129.20$, $p = 0.030$ at lag~1;
$+205.34$, $p = 0.004$ at lag~3. The exit mechanism from
\cite{capponi2024optimal} is confirmed in the high-concentration regime.

\subsection{Test 2: Dose-Response Monotonicity}

\textbf{Hypothesis 2.} Exit rates increase with $\Delta$ in the upper
quartiles.

\textbf{Result: PASS (inverted-U).} Q1--Q2 stable (17.8\%), Q3 peak (20.6\%),
Q4 drops (14.4\%). The non-monotonicity motivates the quadratic and is
captured by the model.

\subsection{Test 3: Lag Robustness}

\textbf{Hypothesis 3.} Both $\beta_1$ and $\beta_2$ are significant for at
least 3 of 5 lag windows.

\textbf{Result: PASS.} Lags 1, 2, 3 all have significant quadratic terms
($p < 0.07$). The turning point is stable at 0.08--0.10. Signal dissipates at
lags 5--7, consistent with short-memory PLP reaction.

\subsection{Test 4: IL Orthogonality}

\textbf{Hypothesis 4.} $\beta_1, \beta_2$ persist with and without the IL
control.

\textbf{Result: PASS.} IL coefficient ($\beta_3 = 13.52$) is positive
(higher IL $\to$ more exits) but does not absorb the concentration effect.
The quadratic structure is present in the $\AT$-only model as well.

\subsection{Test 5: JIT Confounding Control}

\textbf{Hypothesis 5.} The concentration effect is not an artifact of JIT
activity.

\textbf{Result: PASS.} Adding a JIT proxy (positions with blocklife $< 100$
blocks) as a control:
\begin{itemize}
  \item JIT count has independent positive effect on exits
    ($\beta_{\text{JIT}} = 0.177$, $p = 0.031$)
  \item The $\AT$ effect becomes \emph{more} significant
    ($p: 0.072 \to 0.022$) with JIT controlled
  \item Confirms: $\AT$ captures concentration \emph{beyond} JIT
\end{itemize}

\subsection{Test 6: Alternative Treatment Timing}

Three treatment constructions tested:

\begin{center}
\begin{tabular}{@{}llll@{}}
\toprule
Treatment & $\beta_1$ (linear) & $p$ & Finding \\
\midrule
Point A\_T (real, lag 1) & $-3.93$ & 0.072 & Negative (shelter) \\
Lifetime mean A\_T & $-7.60$ & 0.444 & Negative, insignificant \\
Deviation from $1/\PosCount$ & $-5.44$ & 0.109 & Negative (Q3 reveals non-linearity) \\
\textbf{Quadratic deviation} & \multicolumn{2}{l}{\textbf{See Result 3}} & \textbf{Inverted-U confirmed} \\
\bottomrule
\end{tabular}
\end{center}

\section{Robustness: Five Linear Specifications}\label{app:robustness}

Before arriving at the quadratic model, we tested five linear specifications
that progressively refined the treatment variable. All failed to produce a
positive $\beta_1$, but each contributed to understanding the data-generating
process.

\subsection{Specification A: Point Real A\_T}

Treatment $= A_{T,t-\ell}^{\text{real}}$ (single day's concentration).

\begin{center}
\begin{tabular}{@{}ccccc@{}}
\toprule
Lag & $\hat{\beta}_1$ & Cluster SE & $p$ & Sig \\
\midrule
1 & $-3.928$ & 2.186 & 0.072 & * \\
2 & $-8.783$ & 2.227 & 0.0001 & *** \\
3 & $-0.414$ & 1.131 & 0.714 & \\
5 & $-6.523$ & 1.650 & 0.0001 & *** \\
7 & $-5.625$ & 2.261 & 0.013 & ** \\
\bottomrule
\end{tabular}
\end{center}

\textbf{Finding:} Consistently negative. A naive interpretation is that
concentration \emph{protects} LPs. But the proxy A\_T ($x_k = 1/\PosCount$)
showed the same pattern ($\beta_1 = -4.90$, $p = 0.15$), suggesting the level
of $\AT$ confounds pool health (volume) with concentration pressure.

\subsection{Specification B: Lifetime Mean A\_T}

Treatment $=$ mean $A_T^{\text{real}}$ over $[\text{mint}_i, t - \ell]$.
Motivated by the hypothesis that passive LPs react to accumulated exposure,
not point shocks.

\begin{center}
\begin{tabular}{@{}ccccc@{}}
\toprule
Lag & $\hat{\beta}_1$ & Cluster SE & $p$ & Sig \\
\midrule
1 & $-7.595$ & 9.913 & 0.444 & \\
3 & $-11.603$ & 7.959 & 0.145 & \\
5 & $-14.977$ & 8.021 & 0.062 & * \\
7 & $-3.042$ & 6.829 & 0.656 & \\
\bottomrule
\end{tabular}
\end{center}

\textbf{Finding:} Amplifies the negative effect but with much larger SEs
(reduced variation in running means). Marginal significance only at lag~5.

\subsection{Specification C: Deviation from $1/\PosCount$ (Wrong Null)}

First attempt at the deviation treatment, using the old proxy A\_T from a
\emph{different} position set (Q4 burns) as the null. The real and null A\_T
were computed from different populations---an apples-to-oranges comparison.

\textbf{Finding:} Real A\_T was $10$--$100\times$ smaller than the proxy null
(different position sets, different blocklives). Deviation was negative for
40/41 days. Only lag-1 showed positive $\beta_1$ ($+1.92$, $p = 0.43$).
\textbf{Discarded} after identifying the data mismatch.

\subsection{Specification D: Corrected Deviation (Same Position Set)}

Null A\_T recomputed from the \emph{same} Dune query (same positions, same
blocklives). Now 26/41 days have $\Delta > 0$, mean ratio $= 2.65\times$.

\textbf{Finding:} Still negative $\beta_1$ in linear model, but dose-response
revealed the Q3 peak (20.56\%) and Q4 dip (14.37\%) that motivated the
quadratic specification.

\subsection{Specification E: JIT-Controlled Linear}

Added daily JIT count (positions with blocklife $< 100$ blocks) as control.

\textbf{Finding:} $\beta_1$ became \emph{more} negative and more significant
($p: 0.072 \to 0.022$). JIT itself predicts exits but does not confound the
$\AT$ effect.

\section{Dune Queries}\label{app:queries}

\begin{center}
\begin{tabular}{@{}llll@{}}
\toprule
Query & Description & Rows & Cost \\
\midrule
Q4v2 & Position-level burns, blocklife & 600 & 0.18 cr \\
Q5 & Daily IL proxy & 41 & 0.14 cr \\
Q6 & Per-position Collect fees + blocklife (real $\AT$ and null $\AT$ per day)
  & 41 (aggregated) & 3.3 cr \\
\bottomrule
\end{tabular}
\end{center}

Total cost: ${\sim}3.6$ credits.

\textbf{Pool:} ETH/USDC 30bps (\texttt{0x8ad599c3a0ff1de082011efddc58f1908eb6e6d8}).\\
\textbf{Window:} 2025-12-05 to 2026-01-14 (41 days).
